\chapter{Colds}
\index{Colds}

\medskip
\noindent\textit{\textbf{Under review.} This chapter is an AI-assisted educational draft; it carries no source citations and is not a substitute for professional medical advice.}

\section*{Definition and Overview}
The common cold is a mild viral infection of the nose, throat, and upper airways. It causes a runny or blocked nose, sneezing, sore throat, and cough, and is one of the most frequent illnesses in humans. Adults typically catch a cold two to four times a year, and children more often. Colds are usually harmless and clear on their own within about a week to ten days. There is no cure, but symptoms can be eased with self-care, rest, and simple medicines. Antibiotics do not work against colds because they are caused by viruses, not bacteria.

\section*{Epidemiology}
Colds are the most common acute illness worldwide, affecting people of all ages. Children, who are exposed to many viruses at school and childcare, may have six to ten colds or more each year, while adults average around two to four. Colds are more common in autumn and winter in cooler climates, partly because people spend more time indoors and in close contact. Almost everyone will experience many colds over a lifetime. The illness is usually mild, but it is a major cause of time off work and school, and can be more serious in people with asthma or other chest conditions.

\section*{Aetiology and Pathophysiology}
Colds are caused by many different viruses, the most common being rhinoviruses. Other causes include coronaviruses, respiratory syncytial virus, adenoviruses, and parainfluenza viruses. Because there are so many different viruses and strains, having a cold does not protect against future colds.

\begin{itemize}
    \item The virus is spread through droplets in the air when an infected person coughs or sneezes, and by touching contaminated surfaces then touching the eyes, nose, or mouth
    \item The virus infects the lining of the nose and throat, causing inflammation, swelling, and increased mucus production
    \item The body's immune response, rather than the virus alone, produces many of the symptoms
    \item Symptoms usually appear one to three days after infection and peak within a few days
\end{itemize}

\section*{Clinical Features}
\begin{itemize}
    \item Runny or blocked nose with sneezing
    \item Sore or scratchy throat
    \item Cough, which may persist after other symptoms settle
    \item Mild headache, and general feeling of being unwell
    \item Low-grade fever, more common in children
    \item Watery eyes, mild fatigue, and reduced appetite
    \item Clear nasal discharge that may become thicker and coloured over several days, without necessarily indicating a bacterial infection
\end{itemize}

Symptoms usually appear gradually and last for about a week, though a cough can linger for two to three weeks.

\section*{Differential Diagnosis}
\begin{itemize}
    \item Influenza, which tends to come on suddenly with higher fever, body aches, and severe fatigue
    \item COVID-19, which can cause cold-like symptoms and is confirmed with a rapid antigen or PCR test
    \item Allergic rhinitis (hay fever), which causes sneezing and a runny nose without fever and is often seasonal
    \item Strep throat or tonsillitis, bacterial infections that cause a more severe sore throat with fever and swollen glands
    \item Sinusitis, which causes facial pain, pressure, and a blocked nose that persists beyond the usual cold
    \item In infants and young children, bronchiolitis, croup, and pneumonia can begin with cold-like symptoms
    \item Whooping cough, which starts like a cold but progresses to severe coughing spells
\end{itemize}

\section*{When to Seek Medical Care}
Most colds do not need medical attention. See a doctor if:

\begin{itemize}
    \item Symptoms last longer than ten days or seem to be getting worse
    \item You have a high fever, especially a fever lasting more than a few days
    \item You develop shortness of breath, wheezing, or chest pain
    \item You have severe headache, facial pain, or ear pain
    \item A baby under three months has a fever or difficulty feeding
    \item You have an ongoing condition such as asthma, heart disease, or a weakened immune system
\end{itemize}

\textbf{Contact your local emergency services for:} difficulty breathing, severe chest pain, confusion, or if a child is very drowsy, floppy, or struggling to breathe.

\section*{Diagnosis and Investigations}
\begin{itemize}
    \item \textbf{Clinical assessment:} the diagnosis is usually based on the typical symptoms and examination, and tests are rarely needed
    \item \textbf{Examination:} checking the throat, ears, chest, and temperature to look for complications
    \item \textbf{Throat swab:} may be taken if a bacterial infection such as strep throat is suspected
    \item \textbf{Swab for respiratory viruses:} used to test for COVID-19 or influenza when relevant
    \item \textbf{Chest X-ray:} occasionally used if pneumonia is suspected
\end{itemize}

\section*{Management}
\begin{itemize}
    \item \textbf{Rest and fluids:} get plenty of rest and drink enough fluid to stay hydrated
    \item \textbf{Pain relief:} paracetamol or ibuprofen for fever, sore throat, and headache
    \item \textbf{Nasal congestion:} saline nasal sprays or drops, steam, and a humidifier; decongestant sprays are used only briefly
    \item \textbf{Cough:} honey and lemon drinks can soothe a cough in adults and children over one year
    \item \textbf{Medicines with caution:} follow dosing instructions, avoid giving cough and cold medicines to young children, and check for interactions with other medicines
    \item \textbf{Antibiotics:} not effective against colds, and generally only used if a bacterial complication develops
\end{itemize}

\section*{Complications}
\begin{itemize}
    \item Sinusitis, from swelling blocking the sinuses
    \item Middle ear infection, particularly in children
    \item Exacerbation of asthma, triggering wheeze and worsening breathing
    \item Worsening of chronic lung or heart conditions
    \item Bronchitis or pneumonia, more likely in older people, smokers, and those with chronic illness
    \item In infants, difficulty feeding and dehydration
\end{itemize}

\section*{Prognosis and Outcomes}
The common cold usually resolves on its own within seven to ten days, with the cough occasionally lasting longer. Complications are uncommon in healthy people. There is no way to shorten a cold reliably, and treatments aim at comfort rather than cure. Because immunity to the many cold viruses is short-lived, people can expect to catch colds again throughout life, usually with the same self-limiting course. Repeated colds do not indicate a weak immune system for most people.

\section*{Prevention and Self-Care}
\begin{itemize}
    \item Wash hands regularly with soap and water, or use hand sanitiser
    \item Avoid touching the eyes, nose, and mouth with unwashed hands
    \item Stay away from people who are sick, and stay home when you are unwell
    \item Cover coughs and sneezes with a tissue or the elbow, not the hand
    \item Dispose of used tissues promptly and wash hands after
    \item Keep surfaces such as door handles and phones clean
    \item Maintain general health with sleep, exercise, and a balanced diet
    \item Consider annual influenza vaccination, which does not prevent colds but protects against a more serious infection
\end{itemize}
